%% 箱の中の段組：tcolorbox，minipage，figure
\documentclass[twocolumn]{jlreq}
\usepackage[debug]{evenend}
\usepackage[most]{tcolorbox}
\setlength\columnseprule{0.4pt}
\newcommand\para{吾輩は猫である．名前はまだ無い．どこで生れたかとんと見当がつかぬ．何でも薄暗いじめじめした所でニャーニャー泣いていた事だけは記憶している．\par}
\begin{document}
\evenendcolumns{1}
前の段落．\para
\begin{tcolorbox}[title=tcolorboxの中の2段]
箱の最初の段落．
\begin{multicols}{2}
\para\para\para
\end{multicols}
箱の最後の段落．
\end{tcolorbox}
\begin{tcolorbox}[title=3段，段間1字，罫なし]
\begin{evenendmulticols}[columnsep=1\zw,columnseprule=0pt]{3}
\para\para\para\para
\end{evenendmulticols}
\end{tcolorbox}
\begin{tcolorbox}[title=columnbreak]
\begin{multicols}{3}
\para\columnbreak\para\para
\end{multicols}
\end{tcolorbox}
\noindent\begin{minipage}{.6\linewidth}
\begin{multicols}{2}[前置きの見出し]
\para\para
\end{multicols}
\end{minipage}
\evenendcolumns{2}
\para\para\para
\end{document}
